%% ma: L12-B01-0036
%% lop: 12
%% chuong: 1
%% bai: 1
%% dang: 12.1.22
%% loai: TN4
%% muc: TH
%% dap_an: "D"
%% nguon: "HK1-DeOn_1_1.pdf (M10, Đề ôn tập học kì 1 Toán 12, Đề số 2), Phần I Câu 5"
%% kien_thuc: "Bài 1 (tính đơn điệu và cực trị)"
%% trang_thai: cho-duyet
%% ly_do: "Dạng mới đề xuất 12.1.22"
%% ghi_chu: "Hình vẽ lại theo đồ thị y = −x²/(x + 1) (tiệm cận đứng x = −1, tiệm cận xiên y = −x + 1, cực tiểu (−2; 4), cực đại (0; 0)), khớp hình gốc."
\begin{ex}
Cho hàm số $y=f(x)$ có đồ thị như hình bên.
\begin{center}\begin{tikzpicture}[x=0.6cm,y=0.45cm]
\draw[truc] (-5,0)--(3.5,0) node[below] {$x$};
\draw[truc] (0,-4)--(0,8.2) node[left] {$y$};
\node[below right] at (0,0) {$O$};
\draw[phu] (-1,-4)--(-1,8);
\draw[phu] (-4,5)--(3,-2);
\draw[phu] (-2,0)--(-2,4)--(0,4);
\node[below] at (-2,0) {$-2$}; \node[right] at (0,4) {$4$};
\draw[thick,smooth,samples=80,domain=-5:-1.135] plot (\x,{-\x*\x/(\x+1)});
\draw[thick,smooth,samples=80,domain=-0.82:3.3] plot (\x,{-\x*\x/(\x+1)});
\end{tikzpicture}\end{center}
Phát biểu nào sau đây là đúng?
\choice{Tọa độ điểm cực đại là $(-2;4)$}{Tọa độ điểm cực tiểu là $(-2;0)$}{Tọa độ điểm cực đại là $(-2;0)$}{\True Tọa độ điểm cực tiểu là $(-2;4)$}
\loigiai{Trên nhánh trái, đồ thị đi xuống rồi đi lên khi qua $x=-2$, nên $(-2;4)$ là điểm cực tiểu của đồ thị (điểm cực đại là $O(0;0)$).}
\end{ex}
